% P15-U3: exact assumptions for the computation methods in Kass--Raftery Sec. 4.
\chapter{Computing Bayes Factors}
\label{sec:p15-bayes-factor-computation}

A marginal likelihood is an integral, and its numerical approximation is
a further operation with its own error conditions. This chapter develops
the Laplace, Schwarz/BIC, direct Monte Carlo, importance and harmonic-mean
methods described by Kass and Raftery
\cite[Sec.~4, pp.~777--781]{KassRaftery1995BayesFactors}.
The model and data-measure conventions are those of
\cref{def:p15-model-comparison-experiment}. Likelihood normalizing constants
remain part of every expression unless a common cancellation is proved.

The sample-information index \(n\), the parameter dimension \(d\), and the
simulation count \(M\) have different roles. Increasing \(M\) does not repair
an invalid large-\(n\) approximation. Increasing \(n\) does not make an
arbitrary simulation scheme converge to its intended distribution.

\section{A Laplace theorem with explicit local and tail conditions}
\label{sec:p15-laplace-hypotheses}

\begin{hbdefinition}{def:p15-laplace-family}{A uniformly controlled interior mode}
Fix an integer \(d\geq1\). For each \(n\geq1\), let
\(\Theta_n\subseteq\mathbb R^d\) be measurable and contain the ball
\(B(\theta_n,\rho)\), for a fixed \(\rho>0\). Consider
\[
 I_n=\int_{\Theta_n}e^{nh_n(\theta)}a_n(\theta)\,d\theta .
\]
Here \(h_n\) is measurable with values in \(\mathbb R\cup\{-\infty\}\),
and \(a_n\) is finite, nonnegative and measurable. On the stated ball,
\(h_n\) is real and \(C^4\), \(a_n\) is \(C^2\), and the following
constants are independent of \(n\):
\begin{enumerate}
\item \(\nabla h_n(\theta_n)=0\), and
\(J_n=-D^2h_n(\theta_n)\) has all eigenvalues in \([c,C]\), with
\(0<c\leq C<\infty\).
\item The operator norms of \(D^3h_n,D^4h_n,Da_n,D^2a_n\) on the ball
are at most \(A<\infty\), and
\(0<a_-\leq a_n(\theta_n)\leq a_+<\infty\).
\item For every \(v\) with \(\|v\|<\rho\),
\(h_n(\theta_n+v)-h_n(\theta_n)\leq-c_0\|v\|^2\), where \(c_0>0\).
\item For constants \(A_T,\eta>0\),
\[
 \int_{\Theta_n\setminus B(\theta_n,\rho)}
 e^{n(h_n(\theta)-h_n(\theta_n))}a_n(\theta)\,d\theta
 \leq A_Te^{-\eta n}.
\]
\end{enumerate}
The norm is Euclidean, derivative norms are its induced multilinear norms,
and \(d\theta\) is \(d\)-dimensional Lebesgue measure. The last condition
controls all mass outside the local mode, including other modes and tails.
These are deterministic conditions on the sequence being approximated.
\end{hbdefinition}

\begin{hbtheorem}{thm:p15-laplace-relative-error}{proved=annex:p15.laplace-relative-error}{Relative Laplace accuracy under explicit hypotheses}
Under \cref{def:p15-laplace-family},
\[
 \widehat I_n=e^{nh_n(\theta_n)}a_n(\theta_n)
                  (2\pi/n)^{d/2}(\det J_n)^{-1/2}
\]
satisfies \(I_n/\widehat I_n=1+O(n^{-1})\).
The implied constant depends only on the displayed uniform bounds and
fixed dimension, not on \(n\). In particular \(I_n\) is positive and finite.
For two such marginal integrals,
\(B_{10}=I_{1,n}/I_{0,n}\) and
\(\widehat B_{10}=\widehat I_{1,n}/\widehat I_{0,n}\) satisfy
\(B_{10}/\widehat B_{10}=1+O(n^{-1})\) and
\(\log B_{10}-\log\widehat B_{10}=O(n^{-1})\).
\end{hbtheorem}

This theorem does not claim the hypotheses for every statistical model or
for dimensions increasing with \(n\). To apply it to random data, one must
establish the stated bounds along the data sequences in question. A label
such as regular model is not a substitute for the local and tail premises.

\section{Posterior-mode and likelihood-mode constructions}
\label{sec:p15-laplace-mode-variants}

For an observed data set write \(L_n(\theta)\) for its full likelihood,
\(\pi(\theta)\) for a proper prior density, and
\(I_n=\int L_n(\theta)\pi(\theta)d\theta\). Logs below are taken only in
the positive local neighborhood of the specified mode.

\begin{hbcorollary}{thm:p15-laplace-mode-variants}{proved=annex:p15.laplace-mode-variants}{Two expansions of the same integral}
For the posterior-mode construction let
\[
 q_n=\log(L_n\pi),\quad \widetilde\theta_n\in\arg\max q_n,\quad
 H_n=-D^2q_n(\widetilde\theta_n).
\]
If \(h_n=q_n/n,a_n=1,\theta_n=\widetilde\theta_n\) satisfy the explicit
Laplace hypotheses, then
\[
 \widehat I_{\rm post}
 =(2\pi)^{d/2}(\det H_n)^{-1/2}
                L_n(\widetilde\theta_n)\pi(\widetilde\theta_n)
\]
has relative error \(O(n^{-1})\).

For the likelihood-mode construction let
\[
 \ell_n=\log L_n,\quad \widehat\theta_n\in\arg\max\ell_n,\quad
 F_n=-D^2\ell_n(\widehat\theta_n).
\]
If \(h_n=\ell_n/n,a_n=\pi,\theta_n=\widehat\theta_n\) satisfy those
hypotheses, then
\[
 \widehat I_{\rm MLE}
 =(2\pi)^{d/2}(\det F_n)^{-1/2}
                  L_n(\widehat\theta_n)\pi(\widehat\theta_n)
\]
also has relative error \(O(n^{-1})\). These are separate applications
with separate hypothesis checks.
\end{hbcorollary}

The first expansion absorbs the prior into the exponent and uses posterior
curvature; the second leaves it as an amplitude and uses likelihood
curvature. A strongly varying prior can make the second less accurate at a
given finite sample size \cite[Secs.~4.1.1--4.1.2]{KassRaftery1995BayesFactors}.

\begin{hblemma}{thm:p15-laplace-curvature-replacement}{proved=annex:p15.laplace-curvature-replacement}{Replacing an information matrix}
Suppose \(F_n=nJ_n\), where \(J_n\) is symmetric with all eigenvalues
at least \(c>0\).
Replace it by \(\widetilde F_n=n\widetilde J_n\), where
\(\widetilde J_n\) is symmetric and
\(\|\widetilde J_n-J_n\|\leq\delta_n<c/2\). Holding the mode, likelihood
and prior ordinate fixed, the changed approximation satisfies
\[
 \left|\log\frac{\widehat I(\widetilde F_n)}
                         {\widehat I(F_n)}\right|
 \leq d\delta_n/c.
\]
Thus a bound \(\delta_n=O(n^{-1/2})\) gives relative error
\(O(n^{-1/2})\), and \(O(n^{-1})\) preserves that smaller error order.
An expected-information matrix may be substituted under such a proved
comparison; the name expected information alone does not supply the bound.
\end{hblemma}

\section{The precise fixed-dimensional BIC consequence}
\label{sec:p15-bic-consequence}

\begin{hbtheorem}{thm:p15-bic-leading-order}{proved=annex:p15.bic-leading-order}{Log evidence and discarded constant-order terms}
Consider two models of fixed dimensions \(d_0,d_1\) satisfying the MLE
Laplace hypotheses with \(F_{k,n}=nJ_{k,n}\) and proper prior densities
\(\pi_k\). Define
\[
 c_{k,n}=\frac{d_k}{2}\log(2\pi)+\log\pi_k(\widehat\theta_{k,n})
                         -\frac12\log\det J_{k,n}.
\]
Then \(c_{k,n}=O(1)\) and
\[
 \log I_{k,n}
 =\ell_{k,n}(\widehat\theta_{k,n})-\frac{d_k}{2}\log n
                +c_{k,n}+O(n^{-1}).
\]
A parameter-free model with positive finite likelihood uses
\(I_{k,n}=L_{k,n}\), \(d_k=0,c_{k,n}=0\).
With
\(\mathrm{BIC}_k=-2\ell_{k,n}(\widehat\theta_{k,n})+d_k\log n\),
\[
 \log B_{10}
 =-\tfrac12(\mathrm{BIC}_1-\mathrm{BIC}_0)
                     +c_{1,n}-c_{0,n}+O(n^{-1}).
\]
Discarding the \(c\)-terms gives an \(O(1)\) error in log evidence,
not a relative \(O(n^{-1})\) approximation to the factor. If
\(c_{1,n}-c_{0,n}\to c_*\), the ratio of the BIC-based factor to
the exact factor tends to \(e^{-c_*}\). A relatively negligible error on
the log scale follows if \(|\log B_{10}|\to\infty\); it is not asserted
when the log factor remains bounded.
\end{hbtheorem}

Before discarding terms, the two-model MLE formula contains
\[
\begin{split}
2\log B_{10}={}&2(\ell_1(\widehat\theta_1)-\ell_0(\widehat\theta_0))
 +(d_1-d_0)\log(2\pi)-\log\det F_1+\log\det F_0\\
 &+2\log\pi_1(\widehat\theta_1)-2\log\pi_0(\widehat\theta_0)+O(n^{-1}).
\end{split}
\]
In particular both prior-log terms have coefficient two. The information
growth \(F_n=nJ_n\), rather than a convenient count of rows in a table,
determines the role of \(n\). Different blockwise information rates, singular
curvature, boundary modes and increasing dimension require other results
\cite[Sec.~4.1.3]{KassRaftery1995BayesFactors}.

\begin{hbtheorem}{thm:p15-laplace-normal-example}{proved=annex:p15.laplace-normal-example}{An exact normal calculation and three failure boundaries}
For \(X_n\mid\theta\sim\operatorname{normal}(\theta,\sigma^2/n)\),
\(\theta\sim\operatorname{normal}(\mu,\tau^2)\), with fixed positive
\(\sigma,\tau\), the posterior-mode Laplace calculation is exact.
Writing \(s^2=\sigma^2/n\), the MLE approximation is \(\pi(x_n)\) and
\[
 \frac{\widehat I_{\rm MLE}}{I_n}
 =\sqrt{1+s^2/\tau^2}\,
   \exp\!\left\{-\frac{s^2(x_n-\mu)^2}{2\tau^2(s^2+\tau^2)}\right\}.
\]
For a bounded sequence \(x_n\), the stated hypotheses hold for both
variants and this ratio is \(1+O(n^{-1})\).

For \(x_n=\mu=0,\sigma^2=1,\tau^2=4\), against the point-null model
\(\theta=0\), exact, posterior-Laplace, MLE-Laplace and BIC factors
\(B_{10}\) are respectively
\[
 (1+4n)^{-1/2},\quad(1+4n)^{-1/2},\quad(2\sqrt n)^{-1},\quad n^{-1/2}.
\]
The BIC-to-exact ratio tends to two.

The following proper-prior likelihood examples show why hypotheses matter.
For a uniform prior on \([0,1]\) and likelihood \(e^{-n\theta^2/2}\),
the boundary-mode integral is asymptotic to one half of
\(\sqrt{2\pi/n}\). For a uniform prior on
\((-2,-1/2)\cup(1/2,2)\) and likelihood
\(e^{-n(|\theta|-1)^2/2}\), the two-mode integral is asymptotic to
\((2/3)\sqrt{2\pi/n}\), twice one mode's contribution.
For a uniform prior on \([-1,1]\) and likelihood \(e^{-n\theta^4}\),
the integral is asymptotic to
\(n^{-1/4}\int_0^\infty e^{-u^4}du\), not a dimension-one
\(n^{-1/2}\) rate.
\end{hbtheorem}

These likelihoods can arise from \(n\) successes in Bernoulli experiments
whose success probabilities are the corresponding exponentials with \(n=1\).
They are not unnormalized parameter priors. In the examples, respectively,
the interior-ball, one-mode tail, and positive-curvature premises fail.

\section{Direct simulation and importance sampling}
\label{sec:p15-evidence-monte-carlo}

Fix observed data and suppress their symbol. Let \(\Pi\) be a proper prior,
\(L(\theta)\) the finite nonnegative full likelihood, and
\(I=\int L\,d\Pi\in(0,\infty)\). Put \(T(d\theta)=L(\theta)\Pi(d\theta)\)
and posterior \(P=T/I\). Simulation count is \(M\), not sample-information
index \(n\).

\begin{hbtheorem}{thm:p15-direct-evidence-mc}{proved=annex:p15.direct-evidence-mc}{Independent prior simulation}
For i.i.d. \(\theta_i\sim\Pi\),
\[
 \widehat I_M=M^{-1}\sum_{i=1}^M L(\theta_i)
\]
is unbiased and converges almost surely to \(I\).
If \(\int L^2d\Pi<\infty\), its variance is
\(\{\int L^2d\Pi-I^2\}/M\). Integrability of \(L\), rather than a
finite second moment, suffices for the strong-law conclusion.
\end{hbtheorem}

\begin{hbtheorem}{thm:p15-importance-evidence}{proved=annex:p15.importance-evidence}{Ordinary and self-normalized importance are different}
Let \(Q\) be a proposal probability and suppose \(T\ll Q\).
For i.i.d. \(\theta_i\sim Q\), choose a finite almost-everywhere version
\(V=dT/dQ\). Then \(M^{-1}\sum_iV(\theta_i)\) is unbiased, converges
almost surely to \(I\), and has variance
\(\{\int V^2dQ-I^2\}/M\) when that second moment is finite.

For the source's self-normalized form assume the stronger
\(\Pi\ll Q\), let \(W=d\Pi/dQ\), and use
\[
 \widetilde I_M=
 \frac{\sum_iW(\theta_i)L(\theta_i)}{\sum_iW(\theta_i)}
\]
when its denominator is positive. It is eventually defined almost surely
and converges to \(I\). It need not be unbiased for finite \(M\).
If both measures have densities, \(W=\pi/q\) on \(q>0\);
support domination is essential.
\end{hbtheorem}

Kass and Raftery's equation (10) is the self-normalized ratio. An ordinary
importance mean divides by \(M\), not by the random sum of weights.
A proposal improving concentration can worsen tail behavior; no claim of
improvement follows merely from calling a proposal an importance distribution
\cite[Sec.~4.2]{KassRaftery1995BayesFactors}.

\section{Posterior harmonic means and their support}
\label{sec:p15-harmonic-evidence}

\begin{hbtheorem}{thm:p15-harmonic-evidence}{proved=annex:p15.harmonic-evidence}{Reciprocal identities and independent-simulation convergence}
Let \(G\) be a probability measure with \(G\ll T\), and set
\(H=dG/dT\), with a finite version on the posterior support. Then
\[
 E_P H=1/I,\qquad
 E_P H^2=\frac1I\int H^2\,dT.
\]
For i.i.d. \(\theta_i\sim P\), the mean
\(A_M=M^{-1}\sum_iH(\theta_i)\) is unbiased for \(1/I\) and
converges to \(1/I\) almost surely. Its inverse, when positive, is
eventually defined and converges to \(I\); it is not generally unbiased.

The classical harmonic mean takes \(G=\Pi\). It requires
\(L>0\) \(\Pi\)-almost everywhere, gives \(H=1/L\), and has finite
reciprocal-weight variance exactly when \(\int L^{-1}d\Pi<\infty\).
If \(\Pi(L>0)=s<1\), the same posterior average of \(1/L\) instead
converges to \(s/I\), so its inverse converges to \(I/s\).
For density \(g\) of \(G\) and \(h=L\pi\) of \(T\),
the generalized weight is \(g/h\), and its second-moment requirement
is \(\int_{\{h>0\}}g^2/h<\infty\), with \(g=0\) almost everywhere
where \(h=0\).
\end{hbtheorem}

The average \(g/(L\pi)\) estimates the reciprocal evidence. Inverting
this average estimates evidence consistently under the stated sampling
assumptions, but inversion does not preserve unbiasedness. The distinction
is necessary when reading the printed equation (12) and its description
\cite[Sec.~4.3]{KassRaftery1995BayesFactors}.

\begin{hbtheorem}{thm:p15-evidence-mc-clt}{proved=annex:p15.evidence-mc-clt}{Normal limits with the necessary second moments}
For the independent sampling schemes above, the following limits hold,
where variance zero denotes point mass at zero.
An ordinary mean of weights \(Z\) with mean \(I\) and variance
\(v<\infty\) satisfies
\(\sqrt M(\widehat I_M-I)\Rightarrow\operatorname{normal}(0,v)\).
For self-normalized importance, if
\(v_{\rm SN}=E_Q[W^2(L-I)^2]<\infty\), then
\[
 \sqrt M(\widetilde I_M-I)\Rightarrow\operatorname{normal}(0,v_{\rm SN}).
\]
For a generalized harmonic mean, if \(v_H=\operatorname{Var}_P(H)<\infty\),
\[
 \sqrt M(A_M^{-1}-I)\Rightarrow\operatorname{normal}(0,I^4v_H).
\]
Assign any fixed finite default when a finite-sample denominator is zero;
those events disappear in probability under the stated strong laws.
No corresponding normal limit is asserted when these moment premises fail.
\end{hbtheorem}

\begin{hbtheorem}{thm:p15-mc-evidence-boundaries}{proved=annex:p15.mc-evidence-boundaries}{Calculated support, bias, variance and dependence failures}
The following examples use proper priors and finite likelihoods.
\begin{enumerate}
\item On two parameter values, let \(\Pi=(\varepsilon,1-\varepsilon)\),
\(0<\varepsilon<1\), and \(L=(1,0)\). Then \(I=\varepsilon\).
Direct simulation has relative variance
\((1-\varepsilon)/(M\varepsilon)\) and probability
\((1-\varepsilon)^M\) of returning zero. The posterior is concentrated
on the first value, so the classical harmonic estimate equals one,
not \(\varepsilon\). Proposal \(Q=(1,0)\) gives a valid zero-variance
ordinary importance estimate \(\varepsilon\), but the attempted
self-normalized estimate equals one.
\item With prior \((1/2,1/2)\), likelihood \((1,3)\), proposal
\((3/4,1/4)\), and \(M=1\), evidence is 2. Ordinary importance is
unbiased with variance \(16/3\); self-normalized importance has expectation
\(3/2\). The classical posterior harmonic estimate with one draw has
expectation \(5/2\), despite its reciprocal being unbiased for \(1/2\).
Likelihood values exceeding one are allowed for data densities.
\item Let \(\Pi=\operatorname{normal}(0,1)\), \(L=1\), and
\(Q=\operatorname{normal}(0,1/4)\). Ordinary importance has full support
and mean one but infinite variance.
\item For likelihood \(L(\theta)=\phi_{\theta,s^2}(x)\) and prior
\(\operatorname{normal}(0,\tau^2)\), the classical harmonic reciprocal
weight has infinite variance whenever \(\tau^2\geq s^2\), although
its expectation is finite and the i.i.d. strong law still applies.
\item For prior \((1/2,1/2)\), likelihood \((1/4,3/4)\), the evidence
is \(1/2\) and posterior is \((1/4,3/4)\). Initialize a chain in that
posterior and keep its state fixed forever. Every marginal is posterior,
but the harmonic estimate is forever \(1/4\) or \(3/4\), never \(1/2\).
\end{enumerate}
\end{hbtheorem}

All simulation convergence statements above use the integrable i.i.d.
strong law, not an unstated MCMC theorem. For dependent output one needs
a theorem giving the required ergodic averages, with its actual hypotheses.
Approximate posterior draws also need an approximation-error argument.
Stationarity alone is insufficient, as the last example proves.

This component does not supply adaptive Gaussian-quadrature error bounds,
the one-Newton-step implementation, prior/posterior mixture iteration,
Laplace--Metropolis mode estimation, Savage--Dickey or latent-data ratios,
bridge-sampling optimization, or general Gibbs/Metropolis convergence.
These are distinct developments mentioned in Section 4; the identities and
theorem here do not silently discharge them.
